% 90-120 minute graduate lecture: long-context modeling, effective context, context compression,
% hierarchical and block-sparse attention, and recurrent memory.
% Narrative thread: measure what a long window really delivers, then shrink what the model reads,
% choose where it looks, and carry a state across segments, and finally compare the three.
\documentclass[10pt, aspectratio=169]{beamer}

\usepackage{tikz}
\usetikzlibrary{positioning, arrows.meta, shapes.geometric, fit, backgrounds, decorations.pathreplacing}

\input{preamble/packages}
\input{preamble/commands}
\input{preamble/beamer_settings}

\title[Long-Context Modeling]{Long-Context Modeling}
\subtitle{Context Compression, Hierarchical Attention, and Recurrent Memory}

\begin{document}

\input{sections/cover}
\input{sections/long-context/toc}
\input{sections/long-context/motivation}
\input{sections/long-context/learning-outcomes}
\input{sections/long-context/effective-context}
\input{sections/long-context/prompt-compression}
\input{sections/long-context/kv-cache-compression}
\input{sections/long-context/hierarchical-attention}
\input{sections/long-context/block-sparse-attention}
\input{sections/long-context/segment-recurrence}
\input{sections/long-context/linear-memory}
\input{sections/long-context/choosing-a-technique}
\input{sections/long-context/summary}
\input{sections/long-context/references}

\end{document}
